\section{Las tres capas del presente}

Los resultados empíricos de la sección anterior —en particular la dependencia de los picos hiperestructurales respecto de la reducción previa de antisincronización y su relativa robustez ante ruido— sugieren que la intensificación local de la persistencia no se traduce automáticamente en reorganización estructural. Esto motiva una partición ontológica del presente en tres capas jerárquicas de organización de la persistencia. Cada capa se define por su localización, sus marcadores operacionales y su relación con la antisincronización. Las definiciones se formulan de manera que cada capa superior introduce propiedades relacionales o globales que no se deducen de las capas inferiores tomadas aisladamente.

\subsection{Capa 1: Presente local intensificado}

\begin{itemize}
    \item \textbf{Localización}: Módulo individual (componente, variable o región).
    \item \textbf{Marcadores operacionales}: Hiper-persistencia ($\tau_s$ significativamente superior al promedio reciente del propio módulo, calculado sobre ventanas rodantes) combinada con trapping estructural (LAM $\geq 0.80$ o TT $\geq 3.4$ según umbrales de RQA Config B, o aceleraciones significativas en la serie de $\tau_s$).
    \item \textbf{Características ontológicas}: Estado en el que un módulo exhibe una densidad y adhesión localmente intensificada de su configuración de persistencia. Esta intensificación constituye un modo de ser del presente que no se reduce a valores absolutos de $\tau_s$, sino a una desviación significativa respecto del comportamiento basal del módulo en el mismo régimen.
    \item \textbf{Relación con la antisincronización}: Puede emerger incluso bajo alta antisincronización, pero su capacidad de influir en la organización global del sistema permanece estructuralmente limitada mientras persista la fragmentación (alta desviación de $\tau_s$ entre componentes).
\end{itemize}

\begin{conceptbox}[Definición]
El \textit{presente local intensificado} (Capa 1) es el estado en el que un módulo individual sostiene una configuración de persistencia con \textbf{mayor densidad y adhesión local} de lo característico para ese módulo en el período considerado, identificado operacionalmente por la conjunción de hiper-persistencia (desviación positiva de $\tau_s$) y trapping estructural (RQA Config B).
\end{conceptbox}

\subsection{Capa 2: Coherencia entre presentes locales}

\begin{itemize}
    \item \textbf{Localización}: Relaciones inter-modulares (no localizable en ningún módulo tomado por separado).
    \item \textbf{Marcadores operacionales}: Reducción significativa de la antisincronización $A(k)$ —definida como la desviación estándar de las series $\tau_s$ de los distintos componentes en ventanas temporales rodantes— por debajo de un umbral relativo al baseline del sistema.
    \item \textbf{Características ontológicas}: Introducción de una \textbf{dimensión relacional de coherencia de escalas de persistencia} que no existe en los módulos individuales tomados aisladamente. Los módulos pueden continuar evolucionando con trayectorias distintas, pero lo hacen con ritmos de retención ordinal estadísticamente alineados.
    \item \textbf{Rol funcional y estatus ontológico}: Condición necesaria (aunque no suficiente) para que una intensificación local (Capa 1) pueda propagarse y generar efectos detectables a escala del sistema (Capa 3). La Capa 2 constituye la transición crítica que reduce la fragmentación y aumenta la permeabilidad estructural.
\end{itemize}

La Capa 2 no es un estado de sincronización perfecta ni de identidad de trayectorias dinámicas. Es una propiedad \textit{relacional} (estadística) del acoplamiento: la alineación de las escalas de persistencia $\tau_s$. Esta propiedad relacional es ontológicamente irreducible a cualquier descripción que solo examine las series individuales sin considerar su dispersión mutua.

\textbf{Estatus ontológico de la Capa 2}. La Capa 2 tiene estatus ontológico pleno, pero su naturaleza es primariamente relacional y transicional (un «puente ontológico»). A diferencia de la Capa 1 (intensificación local con densidad propia) y de la Capa 3 (reorganización global del sistema), no añade un «presente» con localización sustantiva equivalente. Su contribución irreducible es la coherencia de escalas como propiedad del acoplamiento; su rol principal es habilitar o bloquear la transición de intensificaciones locales a efectos estructurales. Esta caracterización evita tanto reducirla a un mero estado intermedio sin peso ontológico como atribuirle el mismo tipo de «presente» que las otras capas.

\subsection{Capa 3: Reorganización estructural del sistema}

\begin{itemize}
    \item \textbf{Localización}: El sistema considerado como totalidad (no reductible a la agregación de módulos).
    \item \textbf{Marcadores operacionales}: Cambios detectables en la matriz de coherencias entre componentes, medidos por incrementos significativos en la norma de Frobenius entre matrices de correlación o de transición en ventanas consecutivas; aparición de picos hiperestructurales con adelanto estadísticamente significativo respecto a transiciones inyectadas o detectadas; reorganización global de las relaciones inter-modulares.
    \item \textbf{Características ontológicas}: Emergencia de una \textbf{estructura global de relaciones} que no se reduce a la suma de los presentes locales intensificados ni a la mera alineación de escalas. La norma de Frobenius captura una propiedad del sistema como totalidad (la magnitud del cambio en la geometría relacional) que es irreducible a la descripción de sus partes.
    \item \textbf{Condición de emergencia}: Generalmente requiere la conjunción de intensificación local (Capa 1) y reducción previa de fragmentación (Capa 2). Sin la segunda, la primera tiende a permanecer como «fragmento ontológico» sin alterar la estructura global.
\end{itemize}

\begin{importantbox}
La transición de la Capa 1 a la Capa 3 no es automática. Una intensificación local que ocurre bajo alta antisincronización (alta fragmentación) tiende a permanecer como un «fragmento ontológico»: la probabilidad de que genere reorganización estructural detectable (Capa 3) se reduce sustancialmente. Solo cuando la antisincronización desciende se crean las condiciones de permeabilidad necesarias para que la intensificación local adquiera peso estructural. (Estos efectos de filtrado probabilístico se denominan ocasionalmente, con valor heurístico secundario, «resistencia ascendente» y «resistencia descendente».)
\end{importantbox}

\subsection{Jerarquía y no reductibilidad}

Las tres capas constituyen una jerarquía ontológica estricta: cada capa superior presupone condiciones de la inferior, pero introduce propiedades relacionales o globales nuevas e irreductibles. Se explicita a continuación la propiedad irreducible que añade cada capa superior:

\begin{itemize}
    \item La Capa 2 (coherencia de escalas de persistencia) presupone la existencia de presentes locales (Capa 1), pero añade la propiedad relacional nueva e irreducible de la \textbf{alineación de escalas de persistencia $\tau_s$} (cuantificada por la reducción de la desviación estándar entre componentes). Esta propiedad del acoplamiento no existe en los módulos individuales ni se obtiene de sus $\tau_s$ sin considerar su dispersión mutua.
    \item La Capa 3 (reorganización estructural) presupone Capas 1 y 2, pero introduce la propiedad global nueva e irreducible de la \textbf{reorganización de la matriz de relaciones inter-modulares} (medida por la norma de Frobenius). Esta magnitud de cambio en la geometría relacional global no puede predecirse a partir de intensificaciones locales ni de las escalas individuales.
\end{itemize}

La no reductibilidad se afirma aquí en sentido ontológico débil: las descripciones, predicciones y explicaciones formuladas a nivel de una capa superior requieren variables y métricas (antisincronización como desviación estándar de $\tau_s$, norma de Frobenius de la matriz de coherencias) que no figuran en la caracterización exhaustiva de las capas inferiores. Cada capa superior es ontológicamente dependiente de las inferiores, pero no es epifenoménica respecto de ellas; introduce una forma de organización —relacional en el caso de la Capa 2, estructural-global en el de la Capa 3— a la que las métricas y descripciones de nivel inferior son ciegas por principio. La Capa 2, en particular, tiene el estatus de capa relacional/transicional plena: su propiedad irreducible (coherencia de escalas) es condición de posibilidad para la emergencia de la propiedad de la Capa 3, sin reducirse a un mero «estado» carente de peso ontológico propio.

\begin{figure}[htbp]
\centering
\fbox{\parbox{0.9\textwidth}{\sloppy\small\centering
\textbf{Diagrama conceptual de las tres capas del presente}\\[0.2em]
\textbf{Capa 3} (Reorganización estructural del sistema) $\leftarrow$ requiere baja antisincronización $\leftarrow$ \textbf{Capa 2} (Coherencia entre presentes locales) $\leftarrow$ \textbf{Capa 1} (Presente local intensificado: hiper-persistencia + trapping RQA)\\[0.2em]
\textit{La antisincronización regula la probabilidad de transición entre capas mediante un efecto estructural de filtrado probabilístico (ver texto).}
}}
\caption{Esquema de las tres capas del presente. Cada capa introduce propiedades relacionales o globales nuevas e irreductibles a las capas inferiores (ver texto para la caracterización de la Capa 2 como capa relacional/transicional). La antisincronización actúa como regulador de la permeabilidad entre capas mediante el efecto estructural de filtrado probabilístico.}
\label{fig:capas}
\end{figure}